\subsection{Minimal-group collapse in an actual collage}

All modules, homology, and cohomology in this subsection have rational
coefficients.  The quotient functors and induced connecting sets are actual
finite categories and finite sets.

\begin{lemma}[Proved: an averaged projective column suffices for bridge comparison]
\label{CP2:A-averaged-column-bridge-connectivity}
Let $q:M\twoheadrightarrow H$ be a surjective homomorphism of finite
monoids.  Put $R=\mathbb QM$ and $B=\mathbb QH$.  Suppose there are an
idempotent $p\in R$ and a monoid idempotent $f\in H$ such that
$q(p)=f$ and the restricted quotient map
\[
 Rp\longrightarrow Bf
                                                               \tag{1}
\]
is an isomorphism of left $R$ modules.  For every finite left-$N$,
right-$M$ biset $S$, put $U=S\otimes_M H$.  Then the actual bridge map
\[
 B(N,S,M)\longrightarrow B(N,U,H)
                                                               \tag{2}
\]
is an isomorphism on $H_0,H_1$ and a surjection on $H_2$.
There is no assumption on higher homology.
\end{lemma}
\begin{proof}
This is the projective-column argument of
Lemma~\ref{CP2:A-idempotent-quotient-bridge-connectivity}, with the
algebra idempotent $p$ in place of a monoid idempotent in the source;
we verify the comparison explicitly.  Let $T=\mathbb QN$,
$L=\mathbb Q[S]$, and $L'=\mathbb Q[U]$.  Linearization of the set
coequalizer gives $L'=L\otimes_R B$.  Surjectivity of $q$ makes the
canonical map $S\to U$ surjective: if $q(m)=h$, then
$[s,h]=[sm,1]$.  Thus $|U|\leq|S|$, and the commuting left $N$ action
is preserved.

Take the nonnegative two-sided bar complex
$F=L\otimes_R^{\mathbf L}B$, with its natural augmentation onto
$L'$.  Its kernel $D$ is a complex of left-$T$, right-$B$ bimodules
with $H_i(D)=0$ for $i\leq0$.  The good truncation in positive
homological degrees is natural as a bimodule complex; no left-$T$
splitting is used.

Right multiplication by $f$ is exact.  The left projectivity of $Rp$
and the specific isomorphism (1) give
\[
 Ff\simeq L\otimes_R^{\mathbf L}Bf\simeq Lp,
 \qquad L'f\cong L\otimes_R Bf\cong Lp.
\]
The augmentation $Ff\to L'f$ is the same natural isomorphism on
degree-zero homology, so $Df$ is acyclic.  The augmentation of $Bf$
onto the trivial module has kernel $V$, and its exact sequence
$0\to V\to Bf\to\mathbb Q\to0$ yields
\[
 D\otimes_B^{\mathbf L}\mathbb Q
       \simeq(D\otimes_B^{\mathbf L}V)[1].                    \tag{3}
\]
Since $D$ has homology only in degrees at least one, the right side
has homology only in degrees at least two.  The nonnegative left
$N$ bar preserves this lower degree bound.

Derived associativity identifies the chains of the actual map (2)
with the natural map
\[
 \mathbb Q\otimes_T^{\mathbf L}F\otimes_B^{\mathbf L}\mathbb Q
 \longrightarrow
 \mathbb Q\otimes_T^{\mathbf L}L'\otimes_B^{\mathbf L}\mathbb Q.
\]
Its fiber is the analogous tensor with $D$.  The map is induced by
$q$ on monoid arrows and $s\mapsto[s,1]$ on states, so it is the
actual bridge comparison.  Its fiber has no homology below degree
two by (3).  The resulting long exact sequence proves the claim.
\end{proof}

\begin{theorem}[Proved: rational second cohomology under minimal-group collapse]
\label{CP2:A-minimal-group-collapse-collage-H2}
Let $M$ be a finite monoid with minimal ideal in Rees coordinates
$I\times G\times J$.  Let $q:M\twoheadrightarrow\overline M$ be
the actual quotient which forgets the group coordinate inside this
ideal and leaves all outside elements distinct.  For every finite
left-$N$, right-$M$ biset $S$, the actual collage functor
\[
 \mathcal C(N,S,M)\longrightarrow
 \mathcal C(N,S\otimes_M\overline M,\overline M)
                                                               \tag{4}
\]
induces isomorphisms on rational second homology and cohomology.
Consequently, vanishing of all rational degree-two products on the
target implies their vanishing on the source.  This statement does
not assert preservation of bounded homology of the collage.
\end{theorem}
\begin{proof}
The congruence and the rational homology equivalence
$BM\to B\overline M$ are
Lemma~\ref{A:minimal-group-collapse}.  Choose an idempotent $e$ in
the minimal ideal and identify $eMe$ with its finite group $G$.
Put
\[
 p=\frac1{|G|}\sum_{g\in G}g\in\mathbb QM,
 \qquad f=q(e).
\]
Then $p^2=p$ and $q(p)=f$.  In the minimal-column coordinates
$Me=\{ag:a\in I,\ g\in G\}$, the module $Rp$ has basis
\[
 ap=\frac1{|G|}\sum_{g\in G}ag\qquad(a\in I).
                                                               \tag{5}
\]
Indeed $Rp=Re\,p$, each $agp=ap$, and the displayed basis vectors
have disjoint supports in the monoid basis.  The restricted quotient
map sends $ap$ to the corresponding distinct row element of
$\overline M f$.  It is therefore an isomorphism
$Rp\cong\mathbb Q\overline M f$ of left $R$ modules.  Equivariance
also follows directly because left multiplication permutes the
group coordinate within each image row.  Thus the preceding lemma
applies to the bridge in (4).

In the natural two-block exact sequence
\[
 H_2(E)\longrightarrow H_2(BM)\oplus H_2(BN)
 \longrightarrow H_2(B\mathcal C)
 \longrightarrow H_1(E)
 \longrightarrow H_1(BM)\oplus H_1(BN),                         \tag{6}
\]
the induced bridge map is surjective on the first term and an
isomorphism on the fourth.  The corner map is a rational homology
equivalence on both remaining corner terms.  Thus the maps on the
cokernel and kernel adjacent to the middle term are isomorphisms.
The short exact sequences extracted from (6) imply that the middle
map is an isomorphism too.  Over $\mathbb Q$, its dual is the claimed
cohomology isomorphism.  Since this is induced by the actual functor,
it respects products.  For $S=\varnothing$ the assertion follows
directly from the disjoint union of the two corner comparisons.
\end{proof}

\begin{corollary}[Proved: uniform Rees-ideal source gates opposite $R_3$]
\label{CP2:A-Rees-two-row-right-zero-gates}
The following collages have zero rational degree-two products:
\begin{enumerate}
\item The source minimal ideal has two rows and either two or three
columns, with any finite local group; the opposite corner is $R_3$
and the connecting set has arbitrary finite size.
\item The source minimal ideal has two rows and any number of columns,
with any finite local group; the opposite corner is $R_3$ and there
are at most five connecting states.
\item The source has order at most twelve and its minimal local group
is nontrivial; the opposite corner is $R_3$ and the connecting set
has arbitrary finite size.
\end{enumerate}
\end{corollary}
\begin{proof}
Apply (4).  In the first case the collapsed source has a rectangular
$2\times2$ or $2\times3$ minimal band, so use
Theorems~\ref{CP2:A-two-by-two-band-right-zero-all-states} and
\ref{CP2:A-two-by-three-band-general-right-zero}.  Both target
theorems are unconditional in the number of states and in higher
homology.  In the second case the induced connecting set has at
most five elements, so apply
Theorem~\ref{CP2:A-two-row-band-five-state-right-zero}.  Pullback
through the actual second-cohomology isomorphism proves both claims.

For the third case, a group source is rationally acyclic and gives
only relative degree-two classes, whose products are zero.  For a
proper minimal ideal with one row or one column, the collapsed
source has zero second cohomology by the band--Tor formula; the same
is true of the original source by minimal-group collapse.  Since
$BR_3$ is rationally acyclic, all collage degree-two classes are
again relative to the corners and have zero products.

In the remaining case both row and column counts are at least two.
The cardinality inequality
\[
 |I|\,|G|\,|J|+1\leq |M|\leq12,
 \qquad |G|\geq2
\]
forces $|I|=|J|=2$ and $|G|=2$.  The first part applies.  In
particular this conclusion uses an unconditional target theorem,
not a transfer of boundedness through (4).
\end{proof}
